% \documentclass[12pt,a4paper]{article}
% \usepackage[utf8]{inputenc}
% \usepackage{graphicx}
% \usepackage{cite}
% \usepackage{amsmath}
% \usepackage{amsfonts}
% \usepackage{hyperref}
% \usepackage{caption}
% \usepackage{subcaption}
% \usepackage{float}
% \usepackage{fontspec}

% \title{Resource Access Protocols in Real-Time Embedded Systems}
% \author{Raihan Uddin}
% \date{\today}

% \begin{document}

% \maketitle
% \tableofcontents
% \newpage

\documentclass[conference]{IEEEtran}
% \IEEEoverridecommandlockouts
% The preceding line is only needed to identify funding in the first footnote. If that is unneeded, please comment it out.
%Template version as of 6/27/2024

\usepackage{cite}
\usepackage{amsmath,amssymb,amsfonts}
\usepackage{algorithmic}
\usepackage{graphicx}
\usepackage{textcomp}
\usepackage{xcolor}
\usepackage{newtxtext}
\usepackage{enumitem}
\def\BibTeX{{\rm B\kern-.05em{\sc i\kern-.025em b}\kern-.08em
    T\kern-.1667em\lower.7ex\hbox{E}\kern-.125emX}}
\begin{document}

\title{Resource Access Protocols}
\author{
    \IEEEauthorblockN{Raihan Uddin}
    \IEEEauthorblockA{
        \textit{Electronic Engineering} \\
        \textit{Hochschule Hamm-Lippstadt}\\
        Lippstadt, Germany \\
        raihan.uddin@stud.hshl.de
    }
}
\maketitle

\begin{abstract}
Optimized low-level software code, synthesized from high-level behavioral specifications,
ensures efficient and accurate program execution on hardware by offsetting resource
limitations. As hardware and software are parallelly developed to fulfill system-level goals,
software synthesis can pave the path for early planning of the significant challenge in embedded
field- the hardware constraint; thus, can improve upon performance, cost-efficiency and
flexibility across both domains. C/C++/SystemC are the most common languages used for code
automation. The targeted back-end hardware like MCUs, DSPs, FPGAs and AISPs are
significant beneficials of the synthesis procedures. Wherever design parameters like power
consumption, real-time deadlines, memory allocation are of critical concerns, hardware
software seamless alignment is highly anticipated.
The goal of the seminar paper is to explore some of the core aspects of software synthesis like
scheduling tasks to processors. It will provide a comprehensive understanding of efficient
mapping of software onto embedded processors while adhering to system-level critical
constraints including timing, resource limitations and power consumption.
\end{abstract}



\documentclass{article} 

\begin{document} 

\begin{itemize}
\item Definition of software synthesis
\item Difference between manual coding and automated synthesis
\item Importance in embedded systems (productivity, correctness, optimization)
\end{document}
\section{Y-Chart}
Y-Chart, developed in 1983, also known as the Gajski-Y-Chart, is a design methodology that provides a framework for the design of complex systems, particularly in the context of hardware and software co-design. According to~\cite{Gajski.2009}, the Y-Chart approach is characterized by its three main components or axes: behavior (functionality/specification), design structure (netlist/block diagram), and physical design (layout/board design). These components are represented in a Y-shaped diagram, where each arm of the Y represents a different aspect of the design process.
\begin{figure}[htbp]
    \centering
    \includegraphics[width=0.9\columnwidth]{C:/Users/exone/Documents/GitHub/Hardware_Software_Co-design_seminar/Sections/Diagrams/Screenshot 2025-06-19 202509.png}
    \caption{Y-Chart approach~\cite{Gajski.2009}.}
    \label{fig:Y-Chart}
\end{figure}
The Y-Chart in figure~\ref{fig:Y-Chart} comprises three axes, each representing a different aspect of every design process~\cite{Gajski.2009}. The axes are as follows:
\begin{itemize}
        \item \textbf{Behavior}: ``What'' the system should perform
        \item \textbf{Structure}:  ``How'' the system is organized
        \item \textbf{Layout}: ``Where'' the system is implemented
\end{itemize}
The concentrated circles around the origin of the Y-Chart represent design on different abstraction levels and the name of the each abstraction level originates from the types of the components generated on that abstraction level~\cite{Gajski.2009}. The four levels of abstraction are as following:
\begin{itemize}
    \item \textbf{Circuit level}: The components generated in this level are basic building blocks of the circuit i.e., N-type or P-type transistors.
    \item \textbf{Logic level}: This level utilizes logic gates and flip-flops to generate register-transfer components e.g., registers, register files and also functional units such as ALUs and multipliers~\cite{Gajski.2009}.
    \item \textbf{Processor level}: At this level, standard and custom processors along with special-hardware components e.g., memory controllers, arbiters, bridges, routers, and various interface components are generated~\cite{Gajski.2009}.
    \item \textbf{System level}: At this level, embedded systems are designed using processors, memories, buses and other processor components~\cite{Gajski.2009}, generated on the lower levels.
\end{itemize}
The process of implementing a given behavior into a virtual structure on each abstraction level is known as synthesis. Each abstraction level should also contain a database of components that are comprising the structure. After synthesizing at one level, one can move to next lower abstraction level~\cite{Gajski.2009} and synthesize the structure further. However, if there is predefined library of components available with their structures and physical dimensions, one can proceed with the physical designing that consists of floorplanning, placement and routing on the chip or PC board~\cite{Gajski.2009}. Because of the modern day advancement in available software tools, it is possilbe to start designing or synthesizing directly from the system and processor levels without considering the lower levels. As stated in~\cite{Gajski.2009}, ``every system component except standard processors and memories is synthesized to the logic level, before the physical design is performed on the logic level''. This means that all custom components (e.g., controllers, accelerators etc.) except few are readily available to synthesize to the next lower level of abstraction that is logic level in the next step. These exceptional components- processors and memories are reused directly as IPs. In this regard, it is mentioned in~\cite{Gajski.2009} that ``some components on the processor-and-system levels may be obtained as IPs and not synthesized. Therefore, their structure and physical design are known, at least partially, on the level higher than logic level. In that case, the physical design then may contain components of different sizes and from different levels of abstraction''.


\section{Synthesis at Processor Levels}
\section{Importance of Software Synthesis at Processor Level}
Software synthesis produces binary codes for programmable processing elements~\cite{Gajski.2009}. As shown in Figure~\ref{fig:Software_Synthesis}, a system model of a processor describes the components and the connectivity, whereas the bottom portion of the design flow takes the system model as an input and generates an implementation of each component~\cite{Gajski.2009}.
\begin{figure}[htbp]
    \centering
    \includegraphics[width=0.9\columnwidth]{C:/Users/exone/Documents/GitHub/Hardware_Software_Co-design_seminar/Sections/Diagrams/Software_Synthesis.png}
    \caption{Software Synthesis~\cite{Gajski.2009}.}
    \label{fig:Software_Synthesis}
\end{figure}  
Software development for embedded system covers most of the expeses of producing modern multi-processor systems in compared to hardware development. As there are more capabilities of any processor these days, they need firmware, drivers, OS code among others to work properly. Hence, lies the development of multi-faceted software programs to embed in one processor. However, hand-written codes, like conventional practice, do take lot of time for implementation. Hence, engineers, these days, prefer to implement features of processors in software tools because of their flexibility, reusability and for being easily updatable.


\section{Targeted Languages for Embedded Processor}
There are vast range of languages suitable to program embedded systems, like-- assembly, C, C++, and Java. However, there are differences in many aspects--- ``for example in the abstraction level at which the program is written and with that the granularity of control over the processor''~\cite{Gajski.2009}. Among them, assembly language represent the lowest level of abstraction in software development. That is why it is mentioned by~\cite{Gajski.2009} that the assembly is ``basically a symbolic representation of the processor’s machine code and only minimally abstracts the processors complexity''. It gives control over the processor's internals--- such as controlling the registers and instruction-level operations. That is why, it is also valuable for writing performance-critical or low-level routines. However, assembly is very much dependent on the processor's Instruction Set Architecture (ISA), meaning that each processor family (e.g., ARM, x86, RISC-V) differs in its own unique assembly language.
\section{Importance of RTOS}
The languages, those are mentioned before, are unable to provide supprot for concurrency--- which is integral part of embedded systems~\cite{Gajski.2009}. There lies the imoportance of RTOS existence in any embedded system. According to~\cite{Gajski.2009}, embedded applications are often executed on top of an RTOS which allows concurrent task execution. It has, to some extent, similarities with a general purpose operating system. It is a software layer sitting above the bare processor and manages ``concurrent execution of applications and provides various services for communication and synchronization''~\cite{Gajski.2009}.

\section{Software Synthesis Flow}
In the figure~\ref{fig:Software_Synthesis_Flow}, it is shown that the software synthesis takes a high-level model of a system and automatically generates the low-level embedded code, drivers, and binaries needed to run on a processor. It is, thus, eliminating the need for hard-coding, shortening the development time, reducing error-susceptibility, and allowing testing even before the real hardware is built — by using virtual platforms with instruction set simulators.
\begin{figure}[htbp]
    \centering
    \includegraphics[width=0.9\columnwidth]{C:/Users/exone/Documents/GitHub/Hardware_Software_Co-design_seminar/Sections/Diagrams/Software_Synthesis_Flow.png}
    \caption{Software Synthesis Flow~\cite{Gajski.2009}.}
    \label{fig:Software_Synthesis_Flow}
\end{figure}  
\section{TLM}
TLM (Transaction-Level Modeling) is a way to simulate and design systems by describing high-level data transfers between components through busses. However, it utilizes different semantics instead of focusing on low-level signals or wires.
\begin{figure}[htbp]
    \centering
    \includegraphics[width=0.9\columnwidth]{C:/Users/exone/Documents/GitHub/Hardware_Software_Co-design_seminar/Sections/Diagrams/TLM.png}
    \caption{TLM~\cite{Gajski.2009}.}
    \label{fig:TLM}
\end{figure}  
\section{Target Architecture}
Software synthesis typically targets a multi-core embedded platform instead of single processor. These platforms have several general-purpose processors and hardware accelerators, each with its own local memory, programmable interrupt controller (PIC), and timer for time-based tasks.

Processors are connected to the system through a processor bus, enabling communication with local components, memories etc.
\begin{figure}[htbp]
    \centering
    \includegraphics[width=0.9\columnwidth]{C:/Users/exone/Documents/GitHub/Hardware_Software_Co-design_seminar/Sections/Diagrams/target_arch.png}
    \caption{Generic Target Architecture~\cite{Gajski.2009}.}
    \label{fig:target_arch}
\end{figure}  
\section{Code Generation}
Code generation in software synthesis produces sequential code, typically in C for each task that runs on a programmable processor. An exemplary code generation rule is as follows:
\begin{figure}[htbp]
    \centering
    \includegraphics[width=0.9\columnwidth]{C:/Users/exone/Documents/GitHub/Hardware_Software_Co-design_seminar/Sections/Diagrams/code.png}
    \caption{SystemC task specification~\cite{Gajski.2009}.}
    \label{fig:code}
\end{figure}  

\section{Summary}
Overall, software synthesis narrows the gap between high-level system design and low-level executable code, offering faster development cycles, early validation, hardware-software consistency, and scalability across different architectures. It promises to perform more complex architectural design in future with more computing capacity.



\newpage
\bibliographystyle{IEEEtran}  % Change to apa, plain, acm, etc. if required
\bibliography{Reference_list}  % Your .bib file (no .bib extension)

\end{document}
